\subsection{Radius of convergence}\label{sub:radius}
\R{``for which $x$ does $\sum a_n(x-x_0)^n$ converge?''}{
\begin{rcp}
\item Read off the \hd{centre} $x_0$ (in $(x-2)^n$: $x_0=2$) and the \hd{coefficient} $a_n$ = everything without $(x-x_0)^n$.
\item \hd{Radius $R$} --- both formulas always give the same $R$, so pick whichever is easier to evaluate:\\
\hd{ratio} $\frac1R=\lim\left|\frac{a_{n+1}}{a_n}\right|$ when $a_n$ is built from \textbf{factorials or products} (they cancel);\\
\hd{root} \de{Cauchy--Hadamard} $\frac1R=\limsup\sqrt[n]{|a_n|}$ when $n$ sits in an \textbf{exponent}, and the only one that is always valid.\\
Limit $0$ \ra\ $R=\infty$ (converges everywhere); limit $\infty$ \ra\ $R=0$ (only at $x_0$).
\item \warn \hd{Only even/odd powers} ($\sum a_nx^{2n}$)? Substitute $u=x^2$, find the radius $R_u$ in $u$, then $|x|<\sqrt{R_u}$. Reading $R$ straight off $a_n$ gives the wrong interval.
\item Absolute convergence on $|x-x_0|<R$, divergence on $|x-x_0|>R$.
\item \hd{Endpoints $x=x_0\pm R$ one at a time} --- plug in and use \S\ref{sub:seriestree} (usually Leibniz or $p$-series). The two ends can behave differently.
\item Answer as an interval, e.g.\ $[x_0-R,\,x_0+R)$.
\end{rcp}
}
\E{$\sum_{n\ge1}\frac{n^n(x-2)^n}{(2n)!}$}{
$a_n=\frac{n^n}{(2n)!}$; factorial \ra\ ratio test.
\[\left|\frac{a_{n+1}}{a_n}\right|=\frac{(n+1)^{n+1}}{(2n+2)!}\cdot\frac{(2n)!}{n^n}
=\frac{(n+1)\left(1+\frac1n\right)^n}{(2n+1)(2n+2)}\longrightarrow0,\]
since the numerator grows like $e\,n$ and the denominator like $4n^2$. So $\frac1R=0$, $R=\infty$: \textbf{converges for every $x\in\Rl$}.
}
\F{what a power series gives you for free}{
On $|x-x_0|<R$ the sum is $C^\infty$ and you may \textbf{differentiate and integrate termwise}; $R$ is unchanged. Coefficients are unique: $a_n=\frac{f^{(n)}(x_0)}{n!}$.\\
Convergence is \emph{uniform} on every closed $[x_0-r,x_0+r]$ with $r<R$ (Weierstrass $M$-test), \warn not necessarily on the whole open disc.
}

\subsection{Taylor formula and error bounds}\label{sub:taylorerr}
\D{linear / tangent approximation}{
$f(x)\approx f(x_0)+f'(x_0)(x-x_0)$ --- the $n=1$ Taylor polynomial, i.e.\ the tangent line. Error $\le\frac{\max|f''|}2(x-x_0)^2$.\\
This is also the \emph{definition} of differentiability: the tangent is the unique linear function whose error is $o(|x-x_0|)$.
}
\F{Taylor with Lagrange remainder}{
$f(x)=\underbrace{\sum_{k=0}^n\frac{f^{(k)}(x_0)}{k!}(x-x_0)^k}_{T_n(x)}+R_n(x)$, \quad
$R_n(x)=\frac{f^{(n+1)}(\xi)}{(n+1)!}(x-x_0)^{n+1}$ for some $\xi$ between $x_0$ and $x$; also $R_n=o(|x-x_0|^n)$.\\
\hd{Error bound:} $|R_n(x)|\le\frac{M}{(n+1)!}|x-x_0|^{n+1}$ with $M=\max|f^{(n+1)}|$ on the interval between $x_0$ and $x$.
}
\R{approximate a value and bound the error}{
\begin{rcp}
\item Pick $x_0$ near $x$ where $f$ and its derivatives are known exactly ($0$, $1$, $\frac\pi6$, \dots).
\item Write $T_n$ and evaluate.
\item Bound $M=\max|f^{(n+1)}|$ on the interval (for $\sin,\cos$: $M=1$ always; for $e^x$: the right endpoint).
\item If an accuracy $\eps$ is prescribed, solve $\frac{M}{(n+1)!}|x-x_0|^{n+1}<\eps$ for $n$.
\end{rcp}
\hd{Ex.} $e^{0.1}$, $n=3$: $T_3=1+0.1+0.005+\frac{0.001}6\approx1.10517$; $M<1.2$ \ra\ $|R_3|\le\frac{1.2}{24}10^{-4}=5\cdot10^{-6}$.
}

\subsection{Expanding $f$ into its Taylor series \de{Taylorentwicklung}}\label{sub:taylorexp}
\R{``Entwickeln Sie $f$ um $x_0$'' --- route (B) unless the exercise forbids it}{
\begin{rcp}
\item \hd{(A) By derivatives} --- only for a $T_n$ of \emph{fixed, low} degree, or for an $f$ that is in no table. Compute $f,f',\dots,f^{(n)}$, evaluate at $x_0$, insert into $T_n(x)=\sum_{k=0}^n\frac{f^{(k)}(x_0)}{k!}(x-x_0)^k$. Stop at $n$ --- do not hunt for a pattern.
\item \hd{(B) By transforming a known series} (\S\ref{sub:taylorlist}) --- the only sane route to the \emph{full} series with a general term: substitute, factor, differentiate, integrate, multiply (\S\ref{sub:seriesops}).
\item \hd{Centre $x_0\ne0$: substitute $t:=x-x_0$ first}, rewrite $f$ as a function of $t$, expand in $t$, and only then write $t=x-x_0$ back. Expanding in $x$ and ``shifting afterwards'' is wrong.
\item \hd{Radius, then state the domain} $|x-x_0|<R$ (a separate mark): $R=$ distance from $x_0$ to the nearest death of $f$ (pole, $\ln0$, branch), $\infty$ for $e^x,\sin,\cos$; \warn a substitution $u=cx$ \textbf{rescales} it.
\end{rcp}
\hd{The two follow-up questions:} $f^{(n)}(x_0)=n!\,a_n$ --- read the coefficient off the series, never differentiate $17$ times; and $T_n=$ the series cut after $(x-x_0)^n$, remainder bound in \S\ref{sub:taylorerr}.
}
\E{the two centres, side by side}{
\hd{(a) $x_0=0$, $f(x)=\ln(1+2x)$.} Table: $\ln(1+u)=\sum_{n\ge1}\frac{(-1)^{n+1}}nu^n$, $|u|<1$; substitute $u=2x$:
\[\ln(1+2x)=\sum_{n\ge1}\frac{(-1)^{n+1}2^n}{n}x^n=2x-2x^2+\tfrac83x^3-\dots\]
\warn Radius $|2x|<1\iff|x|<\frac12$, \textbf{not} $1$ (distance to the singularity at $-\frac12$). \ $T_3=2x-2x^2+\frac83x^3$; \ $f'''(0)=3!\,a_3=16$.\\
\hd{(b) $x_0=2$, $f(x)=\frac1x$.} Put $t:=x-2$, then pull the constant out to build a $\frac1{1\pm u}$:
\[\frac1x=\frac1{2+t}=\frac12\cdot\frac1{1+\frac t2}=\sum_{n\ge0}\frac{(-1)^n}{2^{n+1}}(x-2)^n,\quad|x-2|<2,\]
$R=2=$ distance from the centre to the pole at $0$. \hd{Every rational $f$}: shift, then force the denominator into $1\pm u$.
}

\subsection{Limits by Taylor: the fastest tool you have}\label{sub:taylorlimit}
\R{limits $x\to0$ by plugging in series}{
\begin{rcp}
\item Look at the denominator's power $x^m$. You need every series to order $m$ --- one order further if the leading terms cancel.
\item Insert the standard series (\S\ref{sub:taylorlist}), expand, cancel.
\item The surviving lowest-order term is the answer; everything else is $o(x^m)$.
\end{rcp}
\hd{$O(x^k)$ means} ``a rest that is at most a constant times $x^k$ near $0$'' --- a bin for everything too small to matter (full rules \S\ref{sub:diffable}). \hd{Bookkeeping:} $\sin x=x+O(x^3)$, so $\frac{\sin x}x=1+O(x^2)$. Carry the $O(x^k)$ through every step and you cannot lose a term.
}
\E{$\displaystyle\lim_{x\to0}\frac{\sin x-x+\frac{x^3}6}{x^5}$}{
$\sin x=x-\frac{x^3}6+\frac{x^5}{120}-\frac{x^7}{5040}+\dots$\\
Numerator $=\frac{x^5}{120}+O(x^7)$ \ra\ limit $=\boxed{\dfrac1{120}}$.\\
l'Hospital would need \textbf{five} differentiations. Don't.
}
\E{more Taylor limits, one line each}{
\begin{bul}
\item $\frac{e^x-1-x}{x^2}=\frac{x^2/2+O(x^3)}{x^2}\to\frac12$.
\item $\frac{\tan x-x}{x^3}\to\frac13$ \ (since $\tan x=x+\frac{x^3}3+O(x^5)$).
\item $\frac1{\sin x}-\frac1x=\frac{x-\sin x}{x\sin x}=\frac{x^3/6}{x^2}\to0$.
\end{bul}
}

\subsection{Building new series from old ones}\label{sub:seriesops}
\R{never recompute derivatives --- transform instead}{
\begin{rcp}
\item \hd{Substitute:} $e^{-x^2}=\sum\frac{(-1)^nx^{2n}}{n!}$; \ $\frac1{1+x^2}=\sum(-1)^nx^{2n}$.
\item \hd{Differentiate termwise:} $\frac1{(1-x)^2}=\sum_{n\ge1}nx^{n-1}$.
\item \hd{Integrate termwise:} $\arctan x=\int\frac{dx}{1+x^2}=\sum\frac{(-1)^nx^{2n+1}}{2n+1}$; \ $\ln(1+x)=\int\frac{dx}{1+x}$.
\item \hd{Multiply:} collect by total degree (Cauchy product). $e^x\sin x=x+x^2+\frac{x^3}3+O(x^4)$.
\item \hd{Divide:} never do long division. Write the denominator as $1-u$ (pull out its constant term first) and use $\frac1{1-u}=1+u+u^2+\dots$, then multiply. \hd{Ex.} $\tan x=\frac{\sin x}{\cos x}$: $\cos x=1-\underbrace{\left(\frac{x^2}2+O(x^4)\right)}_{u}$, so $\frac1{\cos x}=1+\frac{x^2}2+O(x^4)$, and $\tan x=\left(x-\frac{x^3}6\right)\left(1+\frac{x^2}2\right)+O(x^5)=x+\frac{x^3}3+O(x^5)$.
\item \hd{Shift centre:} to expand at $x_0\ne0$, substitute $t=x-x_0$ and expand in $t$.
\end{rcp}
}
